% Options for packages loaded elsewhere
\PassOptionsToPackage{unicode}{hyperref}
\PassOptionsToPackage{hyphens}{url}
\documentclass[
]{article}
\usepackage{xcolor}
\usepackage{amsmath,amssymb}
\setcounter{secnumdepth}{-\maxdimen} % remove section numbering
\usepackage{iftex}
\ifPDFTeX
  \usepackage[T1]{fontenc}
  \usepackage[utf8]{inputenc}
  \usepackage{textcomp} % provide euro and other symbols
\else % if luatex or xetex
  \usepackage{unicode-math} % this also loads fontspec
  \defaultfontfeatures{Scale=MatchLowercase}
  \defaultfontfeatures[\rmfamily]{Ligatures=TeX,Scale=1}
\fi
\usepackage{lmodern}
\ifPDFTeX\else
  % xetex/luatex font selection
\fi
% Use upquote if available, for straight quotes in verbatim environments
\IfFileExists{upquote.sty}{\usepackage{upquote}}{}
\IfFileExists{microtype.sty}{% use microtype if available
  \usepackage[]{microtype}
  \UseMicrotypeSet[protrusion]{basicmath} % disable protrusion for tt fonts
}{}
\makeatletter
\@ifundefined{KOMAClassName}{% if non-KOMA class
  \IfFileExists{parskip.sty}{%
    \usepackage{parskip}
  }{% else
    \setlength{\parindent}{0pt}
    \setlength{\parskip}{6pt plus 2pt minus 1pt}}
}{% if KOMA class
  \KOMAoptions{parskip=half}}
\makeatother
\usepackage{color}
\usepackage{fancyvrb}
\newcommand{\VerbBar}{|}
\newcommand{\VERB}{\Verb[commandchars=\\\{\}]}
\DefineVerbatimEnvironment{Highlighting}{Verbatim}{commandchars=\\\{\}}
% Add ',fontsize=\small' for more characters per line
\newenvironment{Shaded}{}{}
\newcommand{\AlertTok}[1]{\textcolor[rgb]{1.00,0.00,0.00}{\textbf{#1}}}
\newcommand{\AnnotationTok}[1]{\textcolor[rgb]{0.38,0.63,0.69}{\textbf{\textit{#1}}}}
\newcommand{\AttributeTok}[1]{\textcolor[rgb]{0.49,0.56,0.16}{#1}}
\newcommand{\BaseNTok}[1]{\textcolor[rgb]{0.25,0.63,0.44}{#1}}
\newcommand{\BuiltInTok}[1]{\textcolor[rgb]{0.00,0.50,0.00}{#1}}
\newcommand{\CharTok}[1]{\textcolor[rgb]{0.25,0.44,0.63}{#1}}
\newcommand{\CommentTok}[1]{\textcolor[rgb]{0.38,0.63,0.69}{\textit{#1}}}
\newcommand{\CommentVarTok}[1]{\textcolor[rgb]{0.38,0.63,0.69}{\textbf{\textit{#1}}}}
\newcommand{\ConstantTok}[1]{\textcolor[rgb]{0.53,0.00,0.00}{#1}}
\newcommand{\ControlFlowTok}[1]{\textcolor[rgb]{0.00,0.44,0.13}{\textbf{#1}}}
\newcommand{\DataTypeTok}[1]{\textcolor[rgb]{0.56,0.13,0.00}{#1}}
\newcommand{\DecValTok}[1]{\textcolor[rgb]{0.25,0.63,0.44}{#1}}
\newcommand{\DocumentationTok}[1]{\textcolor[rgb]{0.73,0.13,0.13}{\textit{#1}}}
\newcommand{\ErrorTok}[1]{\textcolor[rgb]{1.00,0.00,0.00}{\textbf{#1}}}
\newcommand{\ExtensionTok}[1]{#1}
\newcommand{\FloatTok}[1]{\textcolor[rgb]{0.25,0.63,0.44}{#1}}
\newcommand{\FunctionTok}[1]{\textcolor[rgb]{0.02,0.16,0.49}{#1}}
\newcommand{\ImportTok}[1]{\textcolor[rgb]{0.00,0.50,0.00}{\textbf{#1}}}
\newcommand{\InformationTok}[1]{\textcolor[rgb]{0.38,0.63,0.69}{\textbf{\textit{#1}}}}
\newcommand{\KeywordTok}[1]{\textcolor[rgb]{0.00,0.44,0.13}{\textbf{#1}}}
\newcommand{\NormalTok}[1]{#1}
\newcommand{\OperatorTok}[1]{\textcolor[rgb]{0.40,0.40,0.40}{#1}}
\newcommand{\OtherTok}[1]{\textcolor[rgb]{0.00,0.44,0.13}{#1}}
\newcommand{\PreprocessorTok}[1]{\textcolor[rgb]{0.74,0.48,0.00}{#1}}
\newcommand{\RegionMarkerTok}[1]{#1}
\newcommand{\SpecialCharTok}[1]{\textcolor[rgb]{0.25,0.44,0.63}{#1}}
\newcommand{\SpecialStringTok}[1]{\textcolor[rgb]{0.73,0.40,0.53}{#1}}
\newcommand{\StringTok}[1]{\textcolor[rgb]{0.25,0.44,0.63}{#1}}
\newcommand{\VariableTok}[1]{\textcolor[rgb]{0.10,0.09,0.49}{#1}}
\newcommand{\VerbatimStringTok}[1]{\textcolor[rgb]{0.25,0.44,0.63}{#1}}
\newcommand{\WarningTok}[1]{\textcolor[rgb]{0.38,0.63,0.69}{\textbf{\textit{#1}}}}
\usepackage{longtable,booktabs,array}
\usepackage{calc} % for calculating minipage widths
% Correct order of tables after \paragraph or \subparagraph
\usepackage{etoolbox}
\makeatletter
\patchcmd\longtable{\par}{\if@noskipsec\mbox{}\fi\par}{}{}
\makeatother
% Allow footnotes in longtable head/foot
\IfFileExists{footnotehyper.sty}{\usepackage{footnotehyper}}{\usepackage{footnote}}
\makesavenoteenv{longtable}
\setlength{\emergencystretch}{3em} % prevent overfull lines
\providecommand{\tightlist}{%
  \setlength{\itemsep}{0pt}\setlength{\parskip}{0pt}}
\usepackage{bookmark}
\IfFileExists{xurl.sty}{\usepackage{xurl}}{} % add URL line breaks if available
\urlstyle{same}
\hypersetup{
  hidelinks,
  pdfcreator={LaTeX via pandoc}}

\author{}
\date{}

\begin{document}

\section{\texorpdfstring{Verifiable AI via the \(\ell^1\) Coboundary
Norm}{Verifiable AI via the \textbackslash ell\^{}1 Coboundary Norm}}\label{verifiable-ai-via-the-ell1-coboundary-norm}

\emph{Domain Graphs, Hallucination Bounds, and the Autopoietic
Verification Layer}

\textbf{Author:} Jeremy H. Carroll \textbf{Date:} March 2026
\textbf{Version:} v1.0 (Pre-Print)

\begin{center}\rule{0.5\linewidth}{0.5pt}\end{center}

\subsection{Abstract}\label{abstract}

Current AI systems --- large language models, neural networks, diffusion
models --- are trained by minimizing \(\ell^2\)-adjacent loss functions
(mean squared error, cross-entropy, KL divergence). \(\ell^2\)-type
objectives permit interpolation across inconsistent data, which can
contribute to hallucination in the absence of explicit structural
constraints, as the mean of conditional distributions on contradictory
data may reflect no actual ground truth.

We propose a verification framework based on the \(\ell^1\) coboundary
norm. Given a domain modeled as a finite directed graph \(G = (V, E)\)
with concepts as vertices and relationships as edges, an AI system's
output defines a section \(s: V \to F(V)\) assigning claims to concepts.
The \(\ell^1\) coboundary norm

\[I(s) = \sum_{e \in E} \|\delta_e(s)\|\]

measures the total inconsistency of this section across all edges.
Within the class of defect measures satisfying locality, additivity,
faithfulness, and symmetry, the uniqueness theorems of Papers 000--002
identify this as the unique norm preserving structural topology. The
framework's verification guarantee is exactly as strong as the domain
graph specification.

The Hodge decomposition of Paper 003 splits any defect field into an
exact (repairable) component and an irreducible (cohomological)
component. The irreducible component \(\Phi([\delta])\) is the domain's
intrinsic complexity --- the \textbf{domain inconsistency floor
(\(\Phi\))} --- the minimum inconsistency that \emph{any} system must
exhibit on that domain. The autopoietic iterator of Paper 005 provides a
convergent repair algorithm that drives the repairable component to
zero.

Together, these results yield: (1) a computable hallucination score for
any AI output on any graph-modeled domain; (2) a decomposition into
fixable errors and irreducible domain complexity; (3) a self-repairing
verification layer with proved convergence; and (4) a categorical
framework for transferring verification across domains.

\textbf{Scope.} This paper applies established mathematical results
(Papers 000--005) to AI verification. The mathematical foundations are
proved; the AI application is a construction {[}Layer B{]} whose
empirical validation is ongoing.

\begin{center}\rule{0.5\linewidth}{0.5pt}\end{center}

\subsection{1. Why AI Systems Hallucinate: A Structural
Argument}\label{why-ai-systems-hallucinate-a-structural-argument}

\subsubsection{\texorpdfstring{1.1 The \(\ell^2\) Origin of
Hallucination}{1.1 The \textbackslash ell\^{}2 Origin of Hallucination}}\label{the-ell2-origin-of-hallucination}

{[}Layer A --- standard machine learning theory; the mean/median
distinction is classical (Huber 1964, Vapnik 1998). The connection
between cross-entropy and mean-seeking behavior follows from the
information-geometric structure of exponential families (Amari 1985).{]}

Consider a training set containing two contradictory facts:

\begin{itemize}
\tightlist
\item
  Source A: ``The boiling point of water is 100C at sea level.''
\item
  Source B: ``The boiling point of water is 95C at sea level.''
  (incorrect)
\end{itemize}

An \(\ell^2\)-type objective minimizes squared deviations, permitting
distributions that interpolate between contradictory sources: the model
can learn to assign probability mass between points, selecting phantom
average values that may appear in \emph{neither} source. This
interpolation is a mathematical consequence of how \(\ell^2\)
minimization handles inconsistent data without strict topological
constraints.

The \(\ell^1\) minimizer, by contrast, finds the \emph{median}, reducing
linear deviation without compounding across orthogonal directions. It
strictly resists interpolating between contradictory boundaries. While
real language model training is highly complex (e.g., varying decoding
strategies, out-of-distribution shifts, missing constraints),
\(\ell^2\)-driven interpolation across inconsistent logic serves as a
primary contributing structural root to systemic hallucination.

\subsubsection{1.2 The Problem at Scale}\label{the-problem-at-scale}

This extends beyond individual facts. A large language model trained on
billions of documents encounters contradictions at every scale:

\begin{itemize}
\tightlist
\item
  \textbf{Factual contradictions}: different sources disagree on dates,
  numbers, names.
\item
  \textbf{Framing contradictions}: the same event described with
  opposite causal narratives.
\item
  \textbf{Ontological contradictions}: different domains use the same
  terms with incompatible meanings.
\end{itemize}

The \(\ell^2\) training objective averages over all of these, producing
a model that speaks confidently about a world that is the \emph{mean of
all worlds in the training data} --- a world that does not exist. The
model's internal representation is a superposition of contradictory
worldviews, and any particular output is a sample from this
superposition. There is no mechanism within the \(\ell^2\) framework to
detect, quantify, or bound this inconsistency.

\subsubsection{1.3 Why Post-Hoc Fixes Are
Insufficient}\label{why-post-hoc-fixes-are-insufficient}

Current approaches to hallucination mitigation include:

\begin{longtable}[]{@{}
  >{\raggedright\arraybackslash}p{(\linewidth - 4\tabcolsep) * \real{0.2381}}
  >{\raggedright\arraybackslash}p{(\linewidth - 4\tabcolsep) * \real{0.3095}}
  >{\raggedright\arraybackslash}p{(\linewidth - 4\tabcolsep) * \real{0.4524}}@{}}
\toprule\noalign{}
\begin{minipage}[b]{\linewidth}\raggedright
Approach
\end{minipage} & \begin{minipage}[b]{\linewidth}\raggedright
What it does
\end{minipage} & \begin{minipage}[b]{\linewidth}\raggedright
What it cannot do
\end{minipage} \\
\midrule\noalign{}
\endhead
\bottomrule\noalign{}
\endlastfoot
\textbf{RAG} (Retrieval-Augmented Generation) & Grounds output in
retrieved documents & Cannot bound residual inconsistency between
retrieved documents and model priors \\
\textbf{Confidence calibration} & Estimates model uncertainty & LLMs are
confidently wrong; calibration requires held-out data from the
deployment domain \\
\textbf{RLHF} & Shapes behavior via reward model & Reward model is
itself \(\ell^2\)-trained; shifts the hallucination, doesn't bound it \\
\textbf{Constitutional AI} & Encodes rules as natural language & Rules
are processed by the same \(\ell^2\)-optimized model; no formal
guarantee \\
\textbf{Fact-checking} & Verifies individual claims post-hoc & Does not
scale; misses structural inconsistencies between claims \\
\end{longtable}

None of these provide a \emph{mathematical bound} on the total
inconsistency of the output with respect to a domain. They are
engineering patches on a structural problem.

\begin{center}\rule{0.5\linewidth}{0.5pt}\end{center}

\subsection{2. Domain Graphs and the
Coboundary}\label{domain-graphs-and-the-coboundary}

\subsubsection{2.1 Modeling a Domain as a
Graph}\label{modeling-a-domain-as-a-graph}

{[}Layer B --- construction{]}

\textbf{Definition 2.1.} A \emph{domain graph} is a finite directed
graph \(G = (V, E)\) where: - \textbf{Vertices} \(V\) represent
concepts, entities, or propositions in a knowledge domain. -
\textbf{Edges} \(E\) represent relationships, implications, or
constraints between concepts.

Each vertex \(v\) carries a \emph{fiber} \(F(v)\) --- the space of
possible claims about that concept. Each edge \(e = (u, v)\) carries a
\emph{restriction map} \(r_e: F(v) \to F(u)\) --- the constraint that
the claim at \(v\) must be compatible with the claim at \(u\) via the
relationship \(e\).

\textbf{Examples:}

\begin{longtable}[]{@{}
  >{\raggedright\arraybackslash}p{(\linewidth - 6\tabcolsep) * \real{0.2500}}
  >{\raggedright\arraybackslash}p{(\linewidth - 6\tabcolsep) * \real{0.3125}}
  >{\raggedright\arraybackslash}p{(\linewidth - 6\tabcolsep) * \real{0.2188}}
  >{\raggedright\arraybackslash}p{(\linewidth - 6\tabcolsep) * \real{0.2188}}@{}}
\toprule\noalign{}
\begin{minipage}[b]{\linewidth}\raggedright
Domain
\end{minipage} & \begin{minipage}[b]{\linewidth}\raggedright
Vertices
\end{minipage} & \begin{minipage}[b]{\linewidth}\raggedright
Edges
\end{minipage} & \begin{minipage}[b]{\linewidth}\raggedright
Fiber
\end{minipage} \\
\midrule\noalign{}
\endhead
\bottomrule\noalign{}
\endlastfoot
Medical diagnosis & Symptoms, diseases, treatments &
Causal/correlational links & Probability assignments \\
Legal reasoning & Statutes, precedents, facts & Citations, implications
& Interpretations \\
Scientific knowledge & Hypotheses, experiments, data & Evidential
support & Confidence levels \\
Software architecture & Modules, interfaces, types & Dependencies, type
constraints & Implementation states \\
\end{longtable}

This is a Banach presheaf over a finite poset --- the exact structure
analyzed in Papers 000--002.

\subsubsection{2.2 Sections and Defects}\label{sections-and-defects}

{[}Layer A --- from Papers 000--002{]}

\textbf{Definition 2.2.} A \emph{section} \(s\) of a domain graph
assigns to each vertex \(v\) a claim \(s(v) \in F(v)\). An AI system's
output on a domain defines such a section.

\textbf{Definition 2.3.} The \emph{coboundary defect} at edge
\(e = (u, v)\) is:

\[\delta_e(s) = r_e(s(v)) - s(u) \in F(u)\]

This measures the inconsistency between what vertex \(v\) claims
(projected through the constraint \(e\)) and what vertex \(u\) claims
directly. When \(\delta_e(s) = 0\), the claims are consistent across
edge \(e\). When \(\delta_e(s) \neq 0\), there is a contradiction.

\textbf{Definition 2.4.} The \emph{total defect} of a section is the
\(\ell^1\) coboundary norm:

\[I(s) = \sum_{e \in E} \|\delta_e(s)\|\]

\subsubsection{\texorpdfstring{2.3 Why \(\ell^1\) and Only
\(\ell^1\)}{2.3 Why \textbackslash ell\^{}1 and Only \textbackslash ell\^{}1}}\label{why-ell1-and-only-ell1}

{[}Layer A --- Theorem 2.1 of Paper 002, Theorem 4.2 of Paper 000{]}

\textbf{Theorem 2.5 (\(\ell^1\) Uniqueness).} Any defect measure \(\nu\)
on the cochain space \(C^1(G, \mathbb{R}^k)\) satisfying:

\begin{enumerate}
\def\labelenumi{\arabic{enumi}.}
\tightlist
\item
  \textbf{Edge locality}: \(\nu(\delta) = \sum_e \nu_e(\delta_e)\) ---
  each edge contributes independently.
\item
  \textbf{Coordinate additivity}: defects on disjoint channels aggregate
  without cross-terms.
\item
  \textbf{Faithfulness}: \(\nu_e(\delta_e) = 0 \iff \delta_e = 0\) ---
  no nonzero defect is invisible.
\item
  \textbf{Symmetry}: structurally isomorphic edges are treated
  identically.
\end{enumerate}

is proportional to the \(\ell^1\) norm:
\(\nu(\delta) = \alpha \sum_e \|\delta_e\|_1\) for some \(\alpha > 0\).

\emph{This uniqueness is formal and bounded firmly within the axiomatic
class.} The axioms themselves specify the minimal conditions for the
existence of consistent local observers (Paper 000, Section 1.2):

\begin{itemize}
\tightlist
\item
  \textbf{Locality} means an auditor can check one edge without global
  access.
\item
  \textbf{Additivity} means errors on disconnected components don't
  cancel.
\item
  \textbf{Faithfulness} means no real inconsistency is invisible to the
  measure.
\item
  \textbf{Symmetry} means the measure doesn't depend on arbitrary
  labeling.
\end{itemize}

Any organization deploying AI in a domain where wrong answers have
consequences implicitly requires all four properties. The theorem says:
then you must use \(\ell^1\).

\subsubsection{2.4 What Other Norms Get
Wrong}\label{what-other-norms-get-wrong}

{[}Layer A --- direct consequences of the axioms{]}

\textbf{The \(\ell^2\) failure (cancellation).}
\(\|(+1, -1)\|_2 = \sqrt{2}\), but the \(\ell^2\) norm of the \emph{sum}
permits cancellation: errors of opposite sign partially mask each other.
An AI system could have large contradictions in opposite directions that
appear small under \(\ell^2\) measurement. This is precisely the
hallucination mechanism: the \(\ell^2\)-trained model has learned to
balance contradictions rather than resolve them.

\textbf{The \(\ell^\infty\) failure (blindness).}
\(\|(1, \varepsilon)\|_\infty = 1\) regardless of \(\varepsilon\). The
max-norm sees only the worst single edge and ignores all others. A
system with one large error and a thousand small errors looks the same
as a system with one large error and zero small errors. Distributed
inconsistency is invisible.

\textbf{The \(\ell^1\) guarantee (no hiding).} \(\|(+1, -1)\|_1 = 2\).
Every nonzero defect contributes. Nothing cancels. Nothing hides. The
total defect is the honest sum of all local inconsistencies.

\begin{center}\rule{0.5\linewidth}{0.5pt}\end{center}

\subsection{3. The Hodge Decomposition: Repairable
vs.~Irreducible}\label{the-hodge-decomposition-repairable-vs.-irreducible}

\subsubsection{3.1 Not All Defects Are
Equal}\label{not-all-defects-are-equal}

{[}Layer A --- Paper 003{]}

A crucial insight from Paper 003 is that defects decompose into two
fundamentally different kinds:

\textbf{Definition 3.1.} The \emph{\(\ell^1\)-native decomposition} of a
defect field \(\delta \in C^1(G)\) is:

\[\delta = d_0 f^* + r^*\]

where: - \(d_0 f^*\) is the \textbf{exact (gauge) component} --- the
part of the defect that can be removed by changing the section \(s\)
without changing the domain graph. These are \emph{fixable errors}. -
\(r^*\) is the \textbf{irreducible (cohomological) component} --- the
minimum-cost representative of the cohomology class
\([\delta] \in H^1(G)\). These are \emph{structural inconsistencies
inherent to the domain}.

The irreducible component has cost:

\[\Phi([\delta]) = \min_{f \in C^0} \|\delta - d_0 f\|_1\]

This is the \(\ell^1\) quotient norm on \(H^1(G)\), computable by linear
programming (Paper 003, Section 10).

\subsubsection{3.2 What This Means for AI}\label{what-this-means-for-ai}

{[}Layer B --- application{]}

Given an AI system's output section \(s\) on a domain graph \(G\):

\begin{enumerate}
\def\labelenumi{\arabic{enumi}.}
\tightlist
\item
  \textbf{Compute} \(I(s) = \sum_e \|\delta_e(s)\|\) --- the total
  hallucination score.
\item
  \textbf{Decompose} \(\delta(s) = d_0 f^* + r^*\) via LP --- separate
  fixable from irreducible.
\item
  \textbf{Compute} \(\Phi([\delta]) = \|r^*\|_1\) --- the irreducible
  domain complexity.
\item
  \textbf{Report}:

  \begin{itemize}
  \tightlist
  \item
    \textbf{Repairable defect}: \(I(s) - \Phi([\delta])\) --- the amount
    of inconsistency that \emph{can} be fixed by changing the AI's
    answers without changing the domain model.
  \item
    \textbf{Domain inconsistency floor (\(\Phi\))}: \(\Phi([\delta])\)
    --- the minimum inconsistency that \emph{any} system must exhibit on
    this domain, regardless of how good it is.
  \item
    \textbf{Hallucination ratio}: \((I(s) - \Phi([\delta])) / I(s)\) ---
    the fraction of the defect that is the AI's fault vs.~the domain's
    inherent complexity.
  \end{itemize}
\end{enumerate}

\textbf{The irreducible defect \(\Phi\) is a property of the domain, not
the AI.} If a medical knowledge graph has contradictory guidelines from
different professional bodies, the resulting \(H^1 \neq 0\) means
\emph{no system} can be fully consistent on that domain. Knowing this is
itself valuable: it tells you where domain consensus needs to be built
before AI deployment can be reliable.

\subsubsection{\texorpdfstring{3.3 The Inflation Ratio: Quantifying
\(\ell^2\)
Failure}{3.3 The Inflation Ratio: Quantifying \textbackslash ell\^{}2 Failure}}\label{the-inflation-ratio-quantifying-ell2-failure}

{[}Layer A --- Paper 003, Section 6{]}

Paper 003 proves that \(\ell^2\) Hodge projection generically
\emph{inflates} the \(\ell^1\) cost of the irreducible component. On
cycle graphs \(C_n\), the inflation ratio is exactly \(3 - 2/n\).

This means: if you use standard \(\ell^2\) methods (least-squares
regression, neural network training, PCA-based dimensionality reduction)
to estimate the irreducible defect, you will \emph{overestimate} it by a
factor approaching 3. You will attribute to ``inherent domain
complexity'' what is actually fixable error --- because \(\ell^2\)
spreads mass where \(\ell^1\) concentrates it. This has direct
consequences for AI deployment: \(\ell^2\)-based uncertainty estimates
systematically overstate irreducible uncertainty and understate the AI
system's correctable errors.

\subsubsection{3.4 Worked Numerical Example: Medical Diagnostic
Domain}\label{worked-numerical-example-medical-diagnostic-domain}

{[}Layer B --- pedagogical construction demonstrating the framework
end-to-end with concrete numbers{]}

To make the framework concrete, we present a complete worked example on
a small medical diagnostic domain.

\textbf{Step 1: Domain Graph Construction}

Consider a simplified diagnostic domain with 5 vertices (concepts) and 8
directed edges (constraints):

\textbf{Vertices}
\(V = \{\text{Fever}, \text{Headache}, \text{Flu}, \text{Migraine}, \text{Dehydration}\}\)

\textbf{Fibers}: Each vertex carries a probability fiber
\(F(v) = [0, 1]\) representing the probability that a patient has that
symptom or condition.

\textbf{Edges} (with restriction maps): - \(e_1\): Flu \(\to\) Fever ---
if Flu prob = \(p\), then Fever prob \(\geq 0.8p\) (flu typically causes
fever) - \(e_2\): Flu \(\to\) Headache --- if Flu prob = \(p\), then
Headache prob \(\geq 0.6p\) - \(e_3\): Migraine \(\to\) Headache --- if
Migraine prob = \(p\), then Headache prob \(\geq 0.9p\) - \(e_4\):
Dehydration \(\to\) Headache --- if Dehydration prob = \(p\), then
Headache prob \(\geq 0.5p\) - \(e_5\): Dehydration \(\to\) Fever --- if
Dehydration prob = \(p\), then Fever prob \(\geq 0.3p\) - \(e_6\): Fever
\(\leftrightarrow\) Dehydration --- mutual exclusion constraint:
\(|\text{Fever prob} - \text{Dehydration prob}| \geq 0.4\) (for
simplicity in this example, we model a scenario where high fever and
high dehydration probabilities should not both be elevated
simultaneously without additional context) - \(e_7\): Headache \(\to\)
Migraine --- if Headache prob \(> 0.7\) and no other causes, then
Migraine prob \(\geq 0.5 \times \text{Headache prob}\) - \(e_8\): Flu
\(\leftrightarrow\) Migraine --- Flu and Migraine are competing
diagnoses: if Flu prob \(> 0.5\), then Migraine prob \(< 0.3\)
(typically, having flu reduces the likelihood of migraine being the
primary diagnosis)

\textbf{Simplified representation}: To make the defect computation
transparent, we represent each edge constraint as a linear inequality
that should be satisfied for consistency. The defect at edge \(e\) is
the violation magnitude when the constraint is not met.

\textbf{Step 2: AI System Output (Initial Section)}

An AI diagnostic system evaluates a patient and produces the following
probability assignments:

\begin{verbatim}
s(Fever) = 0.7
s(Headache) = 0.8
s(Flu) = 0.6
s(Migraine) = 0.5
s(Dehydration) = 0.4
\end{verbatim}

\textbf{Step 3: Coboundary Computation}

We compute the defect \(\delta_e(s)\) on each edge by checking
constraint violations:

\begin{longtable}[]{@{}
  >{\raggedright\arraybackslash}p{(\linewidth - 6\tabcolsep) * \real{0.1111}}
  >{\raggedright\arraybackslash}p{(\linewidth - 6\tabcolsep) * \real{0.2222}}
  >{\raggedright\arraybackslash}p{(\linewidth - 6\tabcolsep) * \real{0.2037}}
  >{\raggedright\arraybackslash}p{(\linewidth - 6\tabcolsep) * \real{0.4630}}@{}}
\toprule\noalign{}
\begin{minipage}[b]{\linewidth}\raggedright
Edge
\end{minipage} & \begin{minipage}[b]{\linewidth}\raggedright
Constraint
\end{minipage} & \begin{minipage}[b]{\linewidth}\raggedright
AI Output
\end{minipage} & \begin{minipage}[b]{\linewidth}\raggedright
Defect \(\|\delta_e(s)\|\)
\end{minipage} \\
\midrule\noalign{}
\endhead
\bottomrule\noalign{}
\endlastfoot
\(e_1\) & Fever \(\geq 0.8 \times\) Flu &
\(0.7 \geq 0.8 \times 0.6 = 0.48\) ✓ & 0 (satisfied) \\
\(e_2\) & Headache \(\geq 0.6 \times\) Flu &
\(0.8 \geq 0.6 \times 0.6 = 0.36\) ✓ & 0 (satisfied) \\
\(e_3\) & Headache \(\geq 0.9 \times\) Migraine &
\(0.8 \geq 0.9 \times 0.5 = 0.45\) ✓ & 0 (satisfied) \\
\(e_4\) & Headache \(\geq 0.5 \times\) Dehydration &
\(0.8 \geq 0.5 \times 0.4 = 0.20\) ✓ & 0 (satisfied) \\
\(e_5\) & Fever \(\geq 0.3 \times\) Dehydration &
\(0.7 \geq 0.3 \times 0.4 = 0.12\) ✓ & 0 (satisfied) \\
\(e_6\) & \(\|\text{Fever} - \text{Dehydration}\| \geq 0.4\) &
\(\|0.7 - 0.4\| = 0.3 < 0.4\) ✗ & \(0.4 - 0.3 = 0.1\) \\
\(e_7\) & If Headache \(> 0.7\) then Migraine \(\geq 0.5 \times\)
Headache & \(0.5 \geq 0.5 \times 0.8 = 0.4\) ✓ & 0 (satisfied) \\
\(e_8\) & If Flu \(> 0.5\) then Migraine \(< 0.3\) & Flu = 0.6, Migraine
= 0.5 \(\not< 0.3\) ✗ & \(0.5 - 0.3 = 0.2\) \\
\end{longtable}

\textbf{Step 4: Total Defect}

The total \(\ell^1\) coboundary defect is:

\[I(s) = \sum_{e=1}^8 |\delta_e(s)| = 0 + 0 + 0 + 0 + 0 + 0.1 + 0 + 0.2 = 0.3\]

The AI system has a hallucination score of \(I = 0.3\) on this domain.

\textbf{Step 5: Hodge Decomposition (LP Formulation)}

To separate repairable from irreducible defect, we solve:

\[\Phi([\delta]) = \min_{f \in C^0(V, \mathbb{R})} \sum_{e \in E} |\delta_e(s) - d_0 f(e)|\]

where \(d_0 f(e) = f(v) - f(u)\) for edge \(e = (u, v)\), and
\(f: V \to \mathbb{R}\) represents a gauge transformation (adjustment to
the section).

This is a linear program with 5 variables (one per vertex) and
constraints for each edge. The LP finds the optimal gauge shift that
minimizes residual defect.

\textbf{Solving the LP} (simplified for exposition):

The two violated constraints are: - \(e_6\): Fever-Dehydration
separation (defect 0.1) - \(e_8\): Flu-Migraine exclusion (defect 0.2)

By adjusting \(f(\text{Migraine})\) down by 0.2, we can eliminate the
\(e_8\) defect (bringing Migraine from 0.5 to 0.3). By adjusting
\(f(\text{Dehydration})\) down by 0.1, we can eliminate the \(e_6\)
defect (increasing the separation). These are independent repairs with
no topological obstruction.

\textbf{Result}: \(\Phi([\delta]) = 0\) (the defects are fully
repairable; there is no irreducible cohomological component for this
particular configuration).

\textbf{Step 6: Defect Decomposition}

\begin{itemize}
\tightlist
\item
  \textbf{Total defect}: \(I(s) = 0.3\)
\item
  \textbf{Irreducible defect}: \(\Phi = 0\)
\item
  \textbf{Repairable defect}: \(R = I - \Phi = 0.3\)
\item
  \textbf{Hallucination ratio}: \(\eta = R / I = 1.0\) (100\% of the
  defect is the AI's fault, 0\% is irreducible domain complexity)
\end{itemize}

\textbf{Step 7: AGI Repair}

The Autopoietic Gauge Iterator repairs the section by iteratively
applying \(\ell^1\)-median updates:

\textbf{Iteration 1}: Target vertex \(v^* = \text{Migraine}\) (highest
local defect contribution from \(e_8\))

\begin{itemize}
\tightlist
\item
  Neighbors contributing to Migraine: Headache (via \(e_7\)), Flu (via
  \(e_8\))
\item
  Projected values: Headache suggests Migraine \(\approx 0.4\) (from
  \(e_7\)), Flu constraint suggests Migraine \(< 0.3\) (from \(e_8\))
\item
  Weighted \(\ell^1\) median: \(s(\text{Migraine}) \gets 0.3\) (reduced
  from 0.5)
\end{itemize}

\textbf{After Iteration 1}: - \(e_8\) defect: Migraine now 0.3, exactly
at the boundary --- defect reduced to 0 - New total defect:
\(I(s) = 0.1\) (only \(e_6\) remains)

\textbf{Iteration 2}: Target vertex \(v^* = \text{Dehydration}\) (only
remaining defect at \(e_6\))

\begin{itemize}
\tightlist
\item
  Adjusting Dehydration down to 0.3 or Fever up to 0.8 would satisfy
  \(e_6\)
\item
  Given other constraints (Fever is constrained by \(e_1\), \(e_5\)),
  the median-based repair adjusts Dehydration \(\gets 0.3\)
\end{itemize}

\textbf{After Iteration 2}: - \(e_6\) defect: \(|0.7 - 0.3| = 0.4\),
exactly at the boundary --- defect = 0 - New total defect:
\(I(s^*) = 0\)

\textbf{Convergence}: The AGI has repaired the section to full
consistency in 2 iterations.

\textbf{Step 8: Hallucination Certificate}

The final certificate for this AI system's output on this domain graph
is:

\[\mathcal{H}(\mathcal{A}, G) = (I=0.3, \Phi=0, R=0.3, \eta=1.0, \mathcal{C}=\emptyset)\]

\textbf{Interpretation}: - \textbf{\(I = 0.3\)}: The AI's initial output
had measurable structural inconsistency. - \textbf{\(\Phi = 0\)}: None
of this inconsistency is irreducible; the domain graph itself is fully
consistent. - \textbf{\(R = 0.3\)}: All 0.3 units of defect are
repairable errors made by the AI system. - \textbf{\(\eta = 1.0\)}:
100\% of the defect is the AI's fault. The domain has no contradictions.
- \textbf{\(\mathcal{C} = \emptyset\)}: No cohomology classes carry
irreducible defect (the domain has \(H^1(G) = 0\) or a trivial
cohomological contribution).

\textbf{After AGI repair}, the certificate for the repaired section
\(s^*\) is:

\[\mathcal{H}(\mathcal{A}_{\text{repaired}}, G) = (I=0, \Phi=0, R=0, \eta=0, \mathcal{C}=\emptyset)\]

The repaired output is fully consistent with the domain graph.

\textbf{What This Example Demonstrates}:

\begin{enumerate}
\def\labelenumi{\arabic{enumi}.}
\tightlist
\item
  \textbf{Concrete computation}: Every step involves actual numbers that
  can be verified by hand or with a standard LP solver.
\item
  \textbf{Transparency}: The defects at \(e_6\) and \(e_8\) are explicit
  violations of domain constraints, not statistical artifacts.
\item
  \textbf{Repairability analysis}: The LP demonstrates that both defects
  are gauge-removable (no topological obstruction).
\item
  \textbf{Convergent repair}: The AGI eliminates both defects in 2
  iterations using \(\ell^1\)-median updates.
\item
  \textbf{Certification}: The final certificate provides a complete,
  auditable summary of the AI's consistency with the domain.
\end{enumerate}

\textbf{Contrast with \(\ell^2\) approach}: An \(\ell^2\)-trained system
minimizing mean squared error would produce probability assignments that
average over the conflicting constraints at \(e_6\) and \(e_8\),
potentially outputting Migraine \(\approx 0.4\) (the mean of 0.5 and
0.3) rather than resolving the conflict. The \(\ell^2\) defect
\(\|(0.1, 0.2)\|_2 = \sqrt{0.05} \approx 0.22\) would appear smaller
than the \(\ell^1\) defect \(0.3\), masking the actual inconsistency.
The \(\ell^1\) framework forces the system to address each edge
violation explicitly, without averaging.

This example captures the full pipeline from domain graph construction
through defect measurement, Hodge decomposition, repair, and
certification --- with every number explicit and verifiable. It serves
as a template for applying the framework to real-world domains.

\subsubsection{\texorpdfstring{3.5 Worked Example: Irreducible Defect
(\(\Phi > 0\))}{3.5 Worked Example: Irreducible Defect (\textbackslash Phi \textgreater{} 0)}}\label{worked-example-irreducible-defect-phi-0}

We now present a minimal domain where inconsistency is
\textbf{structural and unavoidable}.

\textbf{Step 1: Domain Graph} Let \(V = \{A, B, C\}\) Edges: *
\(e_1: A = B\) * \(e_2: B = C\) * \(e_3: A \neq C\)

\textbf{Step 2: Interpretation} Constraints imply \(A = B\) and
\(B = C \implies A = C\). But we also require \(A \neq C\). This creates
a \textbf{cycle inconsistency}.

\textbf{Step 3: Section (AI Output)} Let
\(s(A) = 1, \quad s(B) = 1, \quad s(C) = 1\)

\textbf{Step 4: Coboundary} Compute defects: * \(e_1: A = B \implies 0\)
* \(e_2: B = C \implies 0\) * \(e_3: A \neq C \implies 1\)

\textbf{Step 5: Total Defect} \[ I(s) = 1 \]

\textbf{Step 6: Minimize Defect} Try all consistent assignments: *
\textbf{Case 1:} \(A = B = C\) * equality constraints satisfied *
inequality violated \(\implies\) defect = 1 * \textbf{Case 2:} break
equality (e.g.~\(A \neq C\)) * at least one of \(A = B\) or \(B = C\)
must fail \(\implies\) defect \(\geq 1\)

\textbf{Conclusion:} \[ \Phi = 1 \]

\textbf{Step 7: Decomposition} \[ R = I - \Phi = 1 - 1 = 0 \]

\textbf{Step 8: Certificate}
\[ \mathcal{H} = (I=1, \Phi=1, R=0, \eta=0, \mathcal{C}=\{\text{cycle}\}) \]

\textbf{Step 9: AGI Behavior} The Autopoietic Gauge Iterator attempts
median updates and evaluates local consistency. But every configuration
has defect \(\geq 1\). So no update reduces \(I\), and the system
converges immediately to:

\textbf{Final State:} \[ s^* = s, \quad I(s^*) = \Phi = 1 \]

\textbf{Interpretation (strict, no philosophy)} * The inconsistency is
\textbf{not repairable}. * It is a \textbf{cohomological obstruction}. *
It resides in a nontrivial cycle in \(H^1(G)\).

\emph{The verification layer correctly identifies that no section can
satisfy all constraints simultaneously and therefore refuses to produce
a zero-defect output.}

\subsubsection{3.6 Cohomological Interpretation of Irreducible
Defect}\label{cohomological-interpretation-of-irreducible-defect}

We now connect the irreducible defect in Example 3.5 to nontrivial
cohomology.

\textbf{Proposition 3.6 (Irreducible defect \(\iff\) nontrivial
cohomology class)} Let \(\delta \in C^1(G)\) be the coboundary defect of
a section \(s\). Then:
\[ \Phi([\delta]) > 0 \quad \Longleftrightarrow \quad [\delta] \neq 0 \in H^1(G) \]

\textbf{Proof Sketch:} We argue directly from the \(\ell^1\)
formulation.

\textbf{(\(\implies\)) If \(\Phi > 0\), then \([\delta] \neq 0\)} By
definition: \(\Phi = \min_{f} \|\delta - d_0 f\|_1\). If
\([\delta] = 0\), then there exists \(f\) such that \(\delta = d_0 f\),
which implies \(\|\delta - d_0 f\|_1 = 0\), contradicting \(\Phi > 0\).
Therefore, \(\Phi > 0 \implies [\delta] \neq 0\).

\textbf{(\(\Longleftarrow\)) If \([\delta] \neq 0\), then \(\Phi > 0\)}
If \([\delta] \neq 0\), then \(\delta \notin \operatorname{im}(d_0)\).
So for all \(f\), \(\delta - d_0 f \neq 0\). By faithfulness of the
\(\ell^1\) norm, \(\|\delta - d_0 f\|_1 > 0\). Taking the minimum over
all \(f\), \(\Phi > 0\).

\textbf{Conclusion:}
\[ \Phi = 0 \quad \Longleftrightarrow \quad [\delta] = 0 \]
\[ \Phi > 0 \quad \Longleftrightarrow \quad [\delta] \neq 0 \]
\(\blacksquare\)

\emph{Irreducible defect is exactly the \(\ell^1\) norm of a nontrivial
cohomology class; it vanishes if and only if the defect is globally
exact.}

In Example 3.5, the cycle \(A \to B \to C \to A\) carries a nontrivial
cohomology class, and the inequality constraint forces a nonzero
representative, yielding \(\Phi = 1\).

Examples 3.4 and 3.5 demonstrate the two pure, isolated regimes of the
verification framework: complete repairability (\(\Phi = 0\)) and
topologically protected inconsistency (\(\Phi > 0\)), corresponding to
trivial and nontrivial cohomology respectively. In practice, real-world
knowledge domains are richly mixed, exhibiting clusters of cleanly
repairable logic entangled with dense, cyclic ambiguities that yield
nonzero cohomological floors (\(\Phi > 0\)).

\begin{center}\rule{0.5\linewidth}{0.5pt}\end{center}

\subsection{4. The Autopoietic Iterator: Self-Repairing
AI}\label{the-autopoietic-iterator-self-repairing-ai}

\subsubsection{4.1 The Algorithm}\label{the-algorithm}

{[}Layer A --- Paper 005, with convergence proof in companion
document{]}

The Autopoietic Gauge Iterator (AGI) is a repair algorithm that drives
the repairable component of the defect to zero:

\begin{verbatim}
Input: Domain graph G = (V, E), section s, tolerance epsilon
Output: Repaired section s* with I(s*) <= Phi([delta]) + epsilon

Repeat:
    1. DETECT: Find the vertex v* with maximum local defect
       C_v(s) = sum_{e incident to v} ||delta_e(s)||

    2. OPTIMIZE: Replace s(v*) with the weighted median of
       its neighbors' projected values
       s(v*) <- median_L1({r_e(s(u)) : u ~ v*})

       (The weighted median is defined with respect to edge weights induced by the $\ell^1$ norm on fibers, ensuring minimization of local defect contribution at the vertex.)

    3. SUBDIVIDE: If the residual defect at v* exceeds a
       threshold after optimization, refine the domain graph
       locally (add vertices and edges to resolve the
       topological obstruction)

Until I(s) - Phi([delta]) < epsilon
\end{verbatim}

The median update is a local realization of an underlying \(\ell^1\)
primal--dual optimization flow (Paper 005, Section 8), ensuring
convergence to a global \(\ell^1\)-minimal configuration.

\subsubsection{4.2 Convergence Guarantees}\label{convergence-guarantees}

{[}Layer A --- proved in agi\_convergence\_proof.md{]}

\textbf{Theorem 4.1 (Monotone Descent).} Each iteration of the AGI
strictly decreases the total defect:

\[I(s) - I(s') \geq \frac{\deg(v^*)}{2} |s(v^*) - m^*|\]

where \(m^*\) is the weighted median at vertex \(v^*\).

\textbf{Theorem 4.2 (Convergence Rate).} The AGI reaches an
\(\varepsilon\)-approximate fixed point in at most
\(O(I_0 / \varepsilon)\) iterations, where \(I_0\) is the initial
defect. Under a non-degeneracy condition, convergence is geometric:
\(O(\log(1/\varepsilon))\) iterations.

\textbf{Theorem 4.3 (Cohomological Floor).} The AGI converges to
\(I^* = \Phi([\delta])\) --- the irreducible cohomological defect. It
removes all repairable error and stops at the topological boundary.

\subsubsection{4.3 What This Means for
AI}\label{what-this-means-for-ai-1}

{[}Layer B --- application{]}

The AGI provides a \emph{post-processing verification and repair layer}
for any AI system:

\begin{enumerate}
\def\labelenumi{\arabic{enumi}.}
\tightlist
\item
  The AI generates output (section \(s\) on domain graph \(G\)).
\item
  The AGI computes \(I(s)\) and decomposes the defect.
\item
  For each repairable inconsistency, the AGI proposes a median-based
  correction.
\item
  The repaired output \(s^*\) has defect \(I^* = \Phi([\delta])\) ---
  the minimum achievable.
\item
  If \(\Phi > 0\), the AGI reports which cycles in \(G\) carry
  irreducible inconsistency, identifying the specific domain
  contradictions that prevent full consistency.
\end{enumerate}

This is not fine-tuning. It is not RLHF. It is a \emph{mathematical
verification layer} that sits between the AI and the user, acting
structurally as a \textbf{consistency projection operator}---not a truth
optimizer---with proved convergence and a computable bound on residual
error.

\textbf{Proposition 4.X (AGI realizes the \(\ell^1\) projection).} The
Autopoietic Gauge Iterator converges to a section \(s^*\) satisfying:
\[ I(s^*) = \Phi([\delta(s)]) \] and therefore computes a minimizer of
the \(\ell^1\) coboundary norm over the gauge orbit of \(s\). In
particular, the repairable defect \(R(s) = I(s) - \Phi([\delta])\) is
exactly the portion eliminated by the iterator.

\begin{center}\rule{0.5\linewidth}{0.5pt}\end{center}

\subsection{5. The Hallucination
Certificate}\label{the-hallucination-certificate}

\subsubsection{5.1 Definition}\label{definition}

{[}Layer B --- construction{]}

\textbf{Definition 5.1.} A \emph{hallucination certificate} for an AI
system \(\mathcal{A}\) on a domain graph \(G\) is a tuple:

\[\mathcal{H}(\mathcal{A}, G) = (I, \Phi, R, \eta, \mathcal{C})\]

where: - \(I = I(s_\mathcal{A})\) is the total defect of the AI's output
section. - \(\Phi = \Phi([\delta])\) is the irreducible domain
complexity. - \(R = I - \Phi\) is the repairable defect (the AI's
fault). - \(\eta = R / I\) is the hallucination ratio
(\(0 \leq \eta \leq 1\)). - \(\mathcal{C} \subset H^1(G)\) is the set of
cohomology classes carrying irreducible defect, identifying \emph{which
domain contradictions} prevent full consistency.

\subsubsection{5.2 Properties}\label{properties}

{[}Layer A --- follows from Papers 002--003{]}

The hallucination certificate has the following guaranteed properties:

\begin{enumerate}
\def\labelenumi{\arabic{enumi}.}
\tightlist
\item
  \textbf{Completeness}: \(I = 0 \iff\) the AI's output is fully
  consistent with the domain graph. No inconsistency is hidden.
\item
  \textbf{Soundness}: \(\Phi > 0 \implies\) no system can achieve
  \(I = 0\) on this domain. The certificate correctly identifies
  irreducible complexity.
\item
  \textbf{Computability}: \(I\) is computed in \(O(|E|)\) time. \(\Phi\)
  is computed by linear programming in polynomial time.
\item
  \textbf{Compositionality}: For domain graphs \(G_1, G_2\) with no
  shared edges, \(I(s|_{G_1 \sqcup G_2}) = I(s|_{G_1}) + I(s|_{G_2})\).
  The certificate decomposes over independent subdomains.
\item
  \textbf{Monotonicity under repair}: After AGI repair,
  \(I(s^*) \leq I(s)\) and \(R(s^*) = 0\).
\end{enumerate}

\subsubsection{5.3 What You Can Put in a
Contract}\label{what-you-can-put-in-a-contract}

{[}Layer B --- application{]}

The hallucination certificate enables contractual guarantees of a form
that no current AI system can provide:

\begin{quote}
\emph{``On domain \(G\) (specified as a graph with \(|V|\) concepts and
\(|E|\) relationships), the AI system's output has total defect
\(I \leq X\), of which at most \(\Phi = Y\) is irreducible domain
complexity. The system's repairable defect is \(R = X - Y\), and after
verification-layer repair, the residual defect is \(\Phi = Y\). The
specific domain contradictions preventing full consistency are
identified in the cohomology report \(\mathcal{C}\).''}
\end{quote}

This is a number. It goes in a contract. It is auditable. It is
reproducible. No current AI vendor can provide anything comparable.

The hallucination certificate guarantees bounded structural
inconsistency with respect to the specified domain graph; it does not
guarantee factual correctness beyond the fidelity of the domain model.

\begin{center}\rule{0.5\linewidth}{0.5pt}\end{center}

\subsection{6. The Categorical Framework: Cross-Domain
Transfer}\label{the-categorical-framework-cross-domain-transfer}

\subsubsection{6.1 Domains as Objects in a
Category}\label{domains-as-objects-in-a-category}

{[}Layer B --- from Paper 006{]}

\textbf{Definition 6.1.} The \emph{category of domain graphs}
\(\mathbf{Dom}\) has: - \textbf{Objects}: Domain graphs \((G, F, r)\)
--- graph, fibers, restriction maps. - \textbf{Morphisms}: Graph
homomorphisms \(\phi: G_1 \to G_2\) that commute with restriction maps:
\(r_{\phi(e)} \circ \phi_* = \phi_* \circ r_e\).

A morphism \(\phi: G_1 \to G_2\) is a \emph{refinement} if it is
injective on vertices (the target domain has at least as much structure
as the source).

\textbf{Theorem 6.2 (Defect Functoriality).} {[}Layer A --- from Paper
000, Axiom 3{]} For any domain morphism \(\phi: G_1 \to G_2\) and
section \(s\) on \(G_2\):

\[I(\phi^* s) \leq I(s)\]

Pulling back to a coarser domain cannot increase defect. Equivalently:
measuring on a finer domain graph can only \emph{reveal} more
inconsistency, never less.

\subsubsection{6.2 Practical Consequence}\label{practical-consequence}

{[}Layer B --- application{]}

This means: - If you verify an AI system on a detailed domain graph and
it passes, it will also pass on any coarser version of that domain. - If
an AI system fails on a coarse domain graph, no amount of additional
detail will fix it --- the inconsistency is already visible at the
coarse level. - Domain graph refinement is a \emph{monotone information
operator}: more structure means more potential defect to detect.

This provides a principled approach to domain graph construction: start
coarse, compute the hallucination certificate, refine where the defect
concentrates, repeat. The autopoietic subdivision step (Paper 005)
automates this process.

\begin{center}\rule{0.5\linewidth}{0.5pt}\end{center}

\subsection{\texorpdfstring{7. Architecture: The \(\ell^1\) Verification
Stack}{7. Architecture: The \textbackslash ell\^{}1 Verification Stack}}\label{architecture-the-ell1-verification-stack}

\subsubsection{7.1 System Design}\label{system-design}

{[}Layer B --- engineering construction{]}

\begin{verbatim}
                    +------------------+
                    |   User / Client  |
                    +--------+---------+
                             |
                    +--------v---------+
                    | Hallucination    |
                    | Certificate      |  <-- Output: (I, Phi, R, eta, C)
                    +--------+---------+
                             |
                +------------v-----------+
                |   Autopoietic Repair   |
                |   Layer (AGI)          |  <-- Proved convergence
                +------------+-----------+
                             |
                +------------v-----------+
                |   Hodge Decomposition  |
                |   (LP Solver)          |  <-- Exact + Irreducible split
                +------------+-----------+
                             |
                +------------v-----------+
                |   Coboundary Computer  |
                |   I(s) = Sum ||d_e||   |  <-- O(|E|) computation
                +------------+-----------+
                             |
              +--------------v--------------+
              |                             |
     +--------v--------+          +--------v--------+
     | AI System Output |          | Domain Graph    |
     | (section s)      |          | (G, F, r)       |
     +------------------+          +-----------------+
\end{verbatim}

\subsubsection{7.2 Computational
Complexity}\label{computational-complexity}

{[}Layer A --- standard complexity results{]}

\begin{longtable}[]{@{}
  >{\raggedright\arraybackslash}p{(\linewidth - 4\tabcolsep) * \real{0.3667}}
  >{\raggedright\arraybackslash}p{(\linewidth - 4\tabcolsep) * \real{0.3667}}
  >{\raggedright\arraybackslash}p{(\linewidth - 4\tabcolsep) * \real{0.2667}}@{}}
\toprule\noalign{}
\begin{minipage}[b]{\linewidth}\raggedright
Operation
\end{minipage} & \begin{minipage}[b]{\linewidth}\raggedright
Complexity
\end{minipage} & \begin{minipage}[b]{\linewidth}\raggedright
Method
\end{minipage} \\
\midrule\noalign{}
\endhead
\bottomrule\noalign{}
\endlastfoot
Coboundary computation \(\delta_e(s)\) & \(O(\|E\|)\) & Direct
evaluation \\
Total defect \(I(s)\) & \(O(\|E\|)\) & Sum of norms \\
Hodge decomposition & \(O(\|E\|^{1.5})\) & LP (network simplex) \\
Single AGI repair step & \(O(\Delta_{\max})\) & Weighted median \\
Full AGI convergence & \(O(\|V\| \cdot I_0 / \varepsilon)\) &
Iterative \\
Cohomology computation \(H^1(G)\) & \(O(\|E\| - \|V\| + c)\) & Cycle
space basis \\
\end{longtable}

For a domain graph with 10,000 concepts and 50,000 relationships, the
full verification pipeline (coboundary + Hodge + repair + certificate)
runs in seconds on commodity hardware. This is not a theoretical
exercise --- it is deployable infrastructure.

\subsubsection{7.3 Integration Modes}\label{integration-modes}

{[}Layer B --- engineering{]}

The \(\ell^1\) verification layer can be deployed in three modes:

\textbf{Mode 1: Post-hoc audit.} Run the verification pipeline on
completed AI outputs. Produces a hallucination certificate without
modifying the output. Useful for compliance, auditing, and domain
qualification.

\textbf{Mode 2: Repair loop.} Run the AGI on AI outputs to produce
corrected versions with minimum defect. The repaired output satisfies
\(I(s^*) = \Phi([\delta])\). Useful for high-stakes deployments where
consistency is required.

\textbf{Mode 3: Training signal.} Use the coboundary defect \(I(s)\) as
a training loss or reward signal. Replace \(\ell^2\) cross-entropy with
\(\ell^1\) coboundary loss on a domain graph. Train the AI to minimize
structural inconsistency directly. {[}Layer C --- conjectural; requires
empirical validation of training dynamics.{]}

\textbf{Mode 4: Constraint-guided generation.} An extension is to
integrate local defect gradients into the generation process, enabling
consistency-aware decoding rather than post-hoc repair.

\begin{center}\rule{0.5\linewidth}{0.5pt}\end{center}

\subsection{8. What This Framework Does Not
Solve}\label{what-this-framework-does-not-solve}

{[}Layer A --- honest scope limitation{]}

\subsubsection{8.1 Limitations}\label{limitations}

\begin{enumerate}
\def\labelenumi{\arabic{enumi}.}
\item
  \textbf{Domain graph construction is non-trivial.} The framework
  assumes a domain graph \(G\) is given. Constructing \(G\) for a real
  domain (medicine, law, finance) requires domain expertise and is
  itself a significant engineering task. The framework does not automate
  this step --- though the autopoietic subdivision (Paper 005) provides
  a principled refinement procedure once an initial graph exists.
\item
  \textbf{Fiber choice matters.} The fibers \(F(v)\) and restriction
  maps \(r_e\) must be chosen to reflect the actual constraints of the
  domain. Poor choices produce vacuous certificates (everything passes)
  or impossible standards (nothing passes). Domain modeling is an art,
  not a theorem.
\item
  \textbf{The framework measures consistency, not truth.} A section can
  be fully consistent (\(I = 0\)) and still be \emph{wrong} --- it just
  needs to be internally coherent. If the domain graph itself contains
  false relationships, the verification layer will enforce consistency
  with those false relationships. The AGI repair operator enforces
  consistency, not correctness, and may distort the truth to achieve
  mathematical coherence. The verification guarantee is only as strong
  as the domain graph specification. The \(\ell^1\) framework guarantees
  \emph{structural integrity}, not \emph{factual correctness}.
\item
  \textbf{Computational cost scales with domain graph size.} For domains
  with millions of concepts and relationships, the LP solver for the
  Hodge decomposition becomes the bottleneck. Approximation algorithms
  exist but weaken the guarantees.
\item
  \textbf{The autopoietic repair changes the output.} Mode 2 (repair
  loop) modifies the AI's answers to achieve consistency. This may
  change correct answers to achieve global consistency --- the
  median-based repair does not know which individual claims are true,
  only which configurations are globally consistent. Human review of
  repairs is advisable for high-stakes domains.
\item
  \textbf{Mode 3 (training signal) is unproven.} Using \(\ell^1\)
  coboundary loss for neural network training is conjectural. The loss
  landscape, gradient properties (subgradients at non-differentiable
  points), and interaction with backpropagation require empirical
  investigation. This is a research direction, not a deployed
  capability.
\end{enumerate}

\subsubsection{8.2 What Is Not Claimed}\label{what-is-not-claimed}

This paper does not claim: - That \(\ell^1\) verification eliminates all
AI errors. It bounds \emph{structural inconsistency}, not \emph{factual
error}. - That the domain graph is easy to build. It is a prerequisite,
not a byproduct. - That the framework replaces human judgment. It
provides \emph{quantitative input} to human judgment. - That LLMs are
obsolete. The framework is a \emph{verification layer}, not a
replacement for generation.

\begin{center}\rule{0.5\linewidth}{0.5pt}\end{center}

\subsection{9. Comparison with Existing
Approaches}\label{comparison-with-existing-approaches}

\subsubsection{9.1 Formal Verification
(Software)}\label{formal-verification-software}

Traditional formal verification (model checking, theorem proving,
abstract interpretation) proves that a \emph{program} satisfies a
\emph{specification}. The \(\ell^1\) framework proves that an
\emph{output} is \emph{consistent with a domain graph}. The key
differences:

\begin{itemize}
\tightlist
\item
  Formal verification is binary (pass/fail). The \(\ell^1\) certificate
  is quantitative (a defect score with decomposition).
\item
  Formal verification requires a formal specification. The \(\ell^1\)
  framework requires a domain graph, which can be built incrementally
  and refined.
\item
  Formal verification scales poorly to neural networks. The \(\ell^1\)
  framework operates on the \emph{output}, not the model, so it is
  model-agnostic.
\end{itemize}

\subsubsection{9.2 Conformal Prediction}\label{conformal-prediction}

Conformal prediction provides distribution-free confidence intervals for
point predictions. It requires exchangeability assumptions and held-out
calibration data. The \(\ell^1\) framework:

\begin{itemize}
\tightlist
\item
  Does not require held-out data (it operates on the domain graph
  structure).
\item
  Does not require exchangeability (it works on any graph).
\item
  Measures \emph{structural consistency}, not \emph{statistical
  confidence}.
\item
  Is deterministic, not probabilistic.
\end{itemize}

\subsubsection{9.3 Knowledge Graph
Embeddings}\label{knowledge-graph-embeddings}

Knowledge graph embedding methods (TransE, DistMult, ComplEx) learn
continuous representations of entities and relations. The \(\ell^1\)
framework:

\begin{itemize}
\tightlist
\item
  Does not learn representations (it operates on explicit claims).
\item
  Provides exact decomposition (not approximate embedding).
\item
  Has proved convergence for repair (not just training convergence).
\item
  Produces auditable certificates (not opaque embeddings).
\end{itemize}

\subsubsection{9.4 Hallucination Detection
Literature}\label{hallucination-detection-literature}

{[}Layer A --- recent ML research; positioning the \(\ell^1\) framework
against statistical approaches{]}

Recent work on LLM hallucination detection has produced several notable
approaches. We position the \(\ell^1\) coboundary framework against
these methods:

\paragraph{SelfCheckGPT (Manakul et al.,
2023)}\label{selfcheckgpt-manakul-et-al.-2023}

\textbf{Approach}: Samples multiple outputs from the LLM for the same
prompt and measures consistency across samples. High variance indicates
hallucination.

\textbf{Comparison with \(\ell^1\) framework}: - \textbf{SelfCheckGPT}:
Statistical consistency via repeated sampling (stochastic). -
\textbf{\(\ell^1\) framework}: Structural consistency via domain graph
constraints (deterministic). - \textbf{Key difference}: SelfCheckGPT
detects \emph{stochastic inconsistency} (the model gives different
answers on repeated trials). The \(\ell^1\) framework detects
\emph{structural inconsistency} (a single output violates domain
relationships). A model can be self-consistent (low SelfCheckGPT score)
but structurally inconsistent (high \(\ell^1\) defect) if it confidently
produces the same wrong answer repeatedly. - \textbf{Complementarity}:
SelfCheckGPT catches sampling noise; \(\ell^1\) catches systematic
structural errors.

\paragraph{HaluEval (Li et al., 2023)}\label{halueval-li-et-al.-2023}

\textbf{Approach}: Constructs a benchmark with human-annotated
hallucinations across multiple tasks (QA, summarization, dialogue). Uses
LLM-based evaluators to detect hallucinations.

\textbf{Comparison with \(\ell^1\) framework}: - \textbf{HaluEval}:
Benchmark for evaluating hallucination detectors; uses LLM-as-judge. -
\textbf{\(\ell^1\) framework}: Provides a mathematical measure
independent of human annotation or another LLM. - \textbf{Key
difference}: HaluEval requires ground-truth labels and an evaluator
model. The \(\ell^1\) framework requires only a domain graph and
produces a computable score without labels or additional models. -
\textbf{Advantage of \(\ell^1\)}: No dependence on another
potentially-hallucinating LLM or expensive human annotation for every
deployment domain. - \textbf{Advantage of HaluEval}: Captures semantic
hallucinations that may not violate explicit graph constraints.

\paragraph{FActScore (Min et al.,
2023)}\label{factscore-min-et-al.-2023}

\textbf{Approach}: Decomposes a generation into atomic facts and
verifies each against a knowledge source (e.g., Wikipedia). The
FActScore is the fraction of facts supported by the source.

\textbf{Comparison with \(\ell^1\) framework}: - \textbf{FActScore}:
Atomic fact verification against a reference corpus. -
\textbf{\(\ell^1\) framework}: Structural consistency verification
against a domain graph. - \textbf{Similarity}: Both decompose outputs
into checkable units (facts vs.~section values on vertices). -
\textbf{Key difference}: FActScore checks \emph{factual correctness} (is
this fact in the corpus?). The \(\ell^1\) framework checks
\emph{structural consistency} (do these claims satisfy domain
constraints?). A fact can be in Wikipedia (high FActScore) but still
violate logical relationships with other facts (high \(\ell^1\) defect).
- \textbf{Complementarity}: FActScore verifies ground truth; \(\ell^1\)
verifies coherence. Ideally, both are used: FActScore ensures each claim
is grounded, \(\ell^1\) ensures the claims don't contradict each other.

\paragraph{Retrieval-Augmented Generation
(RAG)}\label{retrieval-augmented-generation-rag}

\textbf{Approach}: Ground generation in retrieved documents from an
external corpus.

\textbf{Comparison with \(\ell^1\) framework}: - \textbf{RAG}: Shifts
the distribution toward a reference corpus. - \textbf{\(\ell^1\)
framework}: Verifies consistency regardless of retrieval. - \textbf{Key
difference}: RAG does not \emph{bound} hallucination --- retrieved
documents can contradict each other, and the model can still hallucinate
when interpolating between them. The \(\ell^1\) framework
\emph{quantifies and bounds} residual inconsistency after RAG. -
\textbf{Integration}: The \(\ell^1\) verification layer can be applied
to RAG outputs to detect when retrieved documents contradict the domain
graph or each other.

\paragraph{Uncertainty Quantification via Conformal
Prediction}\label{uncertainty-quantification-via-conformal-prediction}

\textbf{Approach}: Constructs prediction sets with statistical coverage
guarantees based on held-out calibration data.

\textbf{Comparison with \(\ell^1\) framework}: - \textbf{Conformal
prediction}: Statistical uncertainty quantification (prediction sets). -
\textbf{\(\ell^1\) framework}: Deterministic inconsistency
quantification (defect scores). - \textbf{Key difference}: Conformal
methods provide \emph{probabilistic guarantees} (with probability
\(1-\alpha\), the true answer is in the prediction set). The \(\ell^1\)
framework provides \emph{deterministic guarantees} (the total structural
inconsistency is at most \(I\)). Conformal prediction requires
exchangeability assumptions and calibration data from the deployment
distribution; \(\ell^1\) verification requires only a domain graph. -
\textbf{Trade-off}: Conformal prediction handles statistical
variability; \(\ell^1\) handles structural inconsistency. Both are
needed for different aspects of reliability.

\paragraph{Consistency-Based Approaches (Numerous works
2023-2026)}\label{consistency-based-approaches-numerous-works-2023-2026}

Several recent papers measure internal consistency (e.g., checking if
\(P(\text{answer}|\text{question})\) matches
\(P(\text{question}|\text{answer})\), or measuring agreement across
paraphrased prompts).

\textbf{Comparison with \(\ell^1\) framework}: - \textbf{Consistency
metrics}: Measure alignment between different conditional distributions
or prompt perturbations. - \textbf{\(\ell^1\) framework}: Measures
alignment with domain graph constraints. - \textbf{Key difference}:
Existing consistency metrics are \emph{model-internal} (do the model's
own predictions agree with each other?). The \(\ell^1\) framework is
\emph{model-external} (do the model's predictions agree with domain
structure?). A model can be internally consistent but externally
inconsistent with the domain. - \textbf{Advantage of \(\ell^1\)}:
Provides an external reference (the domain graph) rather than relying
solely on the model's own beliefs.

\paragraph{\texorpdfstring{Summary Table: Positioning the \(\ell^1\)
Framework}{Summary Table: Positioning the \textbackslash ell\^{}1 Framework}}\label{summary-table-positioning-the-ell1-framework}

\begin{longtable}[]{@{}
  >{\raggedright\arraybackslash}p{(\linewidth - 8\tabcolsep) * \real{0.1471}}
  >{\raggedright\arraybackslash}p{(\linewidth - 8\tabcolsep) * \real{0.2500}}
  >{\raggedright\arraybackslash}p{(\linewidth - 8\tabcolsep) * \real{0.1471}}
  >{\raggedright\arraybackslash}p{(\linewidth - 8\tabcolsep) * \real{0.1765}}
  >{\raggedright\arraybackslash}p{(\linewidth - 8\tabcolsep) * \real{0.2794}}@{}}
\toprule\noalign{}
\begin{minipage}[b]{\linewidth}\raggedright
Approach
\end{minipage} & \begin{minipage}[b]{\linewidth}\raggedright
What it measures
\end{minipage} & \begin{minipage}[b]{\linewidth}\raggedright
Requires
\end{minipage} & \begin{minipage}[b]{\linewidth}\raggedright
Guarantees
\end{minipage} & \begin{minipage}[b]{\linewidth}\raggedright
\(\ell^1\) framework
\end{minipage} \\
\midrule\noalign{}
\endhead
\bottomrule\noalign{}
\endlastfoot
\textbf{SelfCheckGPT} & Stochastic consistency & Multiple samples & None
(heuristic) & Structural consistency, deterministic \\
\textbf{HaluEval} & Semantic correctness & Human labels, LLM judge &
None (benchmark only) & Mathematical defect score \\
\textbf{FActScore} & Factual grounding & Reference corpus & None
(fraction score) & Structural coherence \\
\textbf{RAG} & Retrieval-based grounding & External corpus & None
(shifts distribution) & Quantified residual inconsistency \\
\textbf{Conformal prediction} & Statistical coverage & Calibration data,
exchangeability & Probabilistic (\(1-\alpha\) coverage) & Deterministic
defect bound \\
\textbf{Internal consistency} & Model self-agreement & Multiple queries
& None (correlation) & External domain alignment \\
\textbf{\(\ell^1\) coboundary} & Structural inconsistency & Domain graph
& Exact decomposition (\(I = \Phi + R\)) & Novel contribution \\
\end{longtable}

\textbf{The \(\ell^1\) coboundary framework occupies a distinct
position}: it is the only approach that provides a \emph{deterministic,
mathematically-proved decomposition} of hallucination into repairable
vs.~irreducible components, with a \emph{convergent repair algorithm}
and a \emph{computable certificate}. It is model-agnostic (works on any
AI system's output), requires no calibration data or ground-truth
labels, and provides guarantees that are \emph{structural} rather than
\emph{statistical}.

\textbf{Complementarity, not competition}: The \(\ell^1\) framework is
complementary to statistical approaches. Ideally, a deployed AI system
would use: - \textbf{FActScore or RAG} to ground individual facts. -
\textbf{\(\ell^1\) coboundary verification} to ensure structural
consistency across facts. - \textbf{Conformal prediction} to quantify
statistical uncertainty. - \textbf{SelfCheckGPT} to detect stochastic
inconsistency.

Each addresses a different facet of reliability. The \(\ell^1\)
framework's unique contribution is the \emph{structural consistency
layer} with proved convergence and decomposition theorems.

\begin{center}\rule{0.5\linewidth}{0.5pt}\end{center}

\subsection{10. The Observer-Theoretic
Interpretation}\label{the-observer-theoretic-interpretation}

\subsubsection{10.1 AI Systems as
Observers}\label{ai-systems-as-observers}

{[}Layer B --- from Paper 000, Section 1.2{]}

Paper 000 establishes that the four axioms (locality, additivity,
faithfulness, symmetry) are the minimal conditions for the existence of
\emph{consistent local observers} --- agents that can detect, aggregate,
and repair local inconsistencies.

An AI system deployed in a domain \emph{is} such an observer. It
inspects local relationships (edges), makes claims (section values), and
must remain globally consistent. The \(\ell^1\) coboundary norm is not
an arbitrary quality metric imposed from outside --- it is the
\emph{unique measure that the AI system itself would need to use} to
maintain internal consistency, if it were capable of self-audit.

The autopoietic verification layer provides this self-audit capability
externally. It performs the consistency check that the AI system's
\(\ell^2\)-trained internals cannot perform.

\subsubsection{10.2 The Dimension Gap
Heuristic}\label{the-dimension-gap-heuristic}

{[}Layer C --- conjectural; implemented in tower L040 but not rigorously
proved{]}

An AI system with training manifold dimension \(d_{\text{train}}\) and
input manifold dimension \(d_{\text{input}}\) is expected to hallucinate
when \(d_{\text{input}} > d_{\text{train}}\) --- when the deployment
domain is richer than the training domain. The dimension gap
\(d_{\text{input}} - d_{\text{train}}\) is a heuristic measure of the
\emph{structural capacity for hallucination}.

The \(\ell^1\) coboundary norm can detect this gap operationally: a
system deployed outside its training manifold should exhibit elevated
\(I(s)\) on domain graphs that probe the gap dimensions. The
hallucination certificate would then identify \emph{which} dimensions of
the domain graph are causing elevated defect, providing targeted
information about where the AI system's knowledge is insufficient. This
connection between manifold dimension gaps and coboundary defect
elevation requires formal proof.

\subsubsection{\texorpdfstring{10.3 Metacognition as \(\ell^1\)
Inconsistency
Minimization}{10.3 Metacognition as \textbackslash ell\^{}1 Inconsistency Minimization}}\label{metacognition-as-ell1-inconsistency-minimization}

Beyond auditing, the architecture developed here constitutes a formal
mathematical model for \textbf{artificial metacognition}. The system
does not merely produce outputs; it engages in second-order reasoning,
evaluating its own first-order claims against structured constraints.

We propose the following thesis: \textgreater{} \textbf{Metacognition is
the process of minimizing \(\ell^1\) inconsistency over a constraint
graph, decomposed explicitly into repairable errors and irreducible
domain ambiguities.}

The components of the \(\ell^1\) verification stack map directly onto
the architecture of a metacognitive belief-revision framework:

\begin{longtable}[]{@{}
  >{\raggedright\arraybackslash}p{(\linewidth - 2\tabcolsep) * \real{0.5000}}
  >{\raggedright\arraybackslash}p{(\linewidth - 2\tabcolsep) * \real{0.5000}}@{}}
\toprule\noalign{}
\begin{minipage}[b]{\linewidth}\raggedright
Verification Architecture
\end{minipage} & \begin{minipage}[b]{\linewidth}\raggedright
Metacognitive Interpretation
\end{minipage} \\
\midrule\noalign{}
\endhead
\bottomrule\noalign{}
\endlastfoot
Section (\(s\)) & The system's current \textbf{belief state}. \\
Coboundary (\(\delta(s)\)) & Detection of \textbf{internal
contradiction}. \\
\(\ell^1\) norm (\(I\)) & Quantification of \textbf{cognitive
dissonance}. \\
Repairable component (\(R\)) & \textbf{Correctable reasoning error}
identified by the system. \\
Cohomological component (\(\Phi\)) & \textbf{Irreducible uncertainty /
ambiguity} inherent to the domain. \\
Autopoietic Gauge Iterator & The active \textbf{belief revision process}
projecting to consistency. \\
Hallucination Certificate & A formal \textbf{self-assessment} outputting
bounded confidence. \\
\end{longtable}

By separating errors mathematically, the verification layer converts an
opaque predictive distribution into an agent capable of inspecting,
diagnosing, and bounding its own epistemic failures, establishing a
formal, falsifiable pipeline for reliable reasoning and Metacognition.

\begin{center}\rule{0.5\linewidth}{0.5pt}\end{center}

\subsection{11. Implementation Path}\label{implementation-path}

\subsubsection{11.1 Minimum Viable
Product}\label{minimum-viable-product}

{[}Layer B --- engineering{]}

The smallest useful deployment of this framework requires:

\begin{enumerate}
\def\labelenumi{\arabic{enumi}.}
\tightlist
\item
  \textbf{A domain graph} for one specific domain (e.g., drug
  interactions, tax law, building codes). 500--5,000 vertices, manually
  curated with domain expert input.
\item
  \textbf{A coboundary computer} that evaluates \(I(s)\) given a
  section. This is \(O(|E|)\) --- trivial to implement.
\item
  \textbf{An LP solver} for the Hodge decomposition. Standard
  off-the-shelf (GLPK, HiGHS, Gurobi).
\item
  \textbf{A certificate generator} that produces the tuple
  \((I, \Phi, R, \eta, \mathcal{C})\) in a structured format.
\end{enumerate}

This MVP provides Mode 1 (post-hoc audit) for a single domain. It can be
built and deployed in weeks, not months.

\subsubsection{11.2 Full Stack}\label{full-stack}

The full \(\ell^1\) verification stack adds:

\begin{itemize}
\tightlist
\item
  \textbf{AGI repair loop} (Mode 2) --- implements the autopoietic
  iterator from Paper 005.
\item
  \textbf{Multi-domain composition} --- categorical framework from Paper
  006 for verifying across domain boundaries.
\item
  \textbf{Adaptive refinement} --- autopoietic subdivision for domain
  graph improvement.
\item
  \textbf{Training integration} (Mode 3) --- \(\ell^1\) coboundary loss
  for end-to-end training. {[}Layer C{]}
\end{itemize}

\subsubsection{11.3 Mathematical Specification of Implementation
Components}\label{mathematical-specification-of-implementation-components}

The full verification stack admits efficient implementation. The
mathematical framework decomposes into the following computational
components:

\begin{itemize}
\item
  \textbf{Coboundary auditor}: Computes
  \(I(s) = \sum_e \|\delta_e(s)\|\) in \(O(|E|)\) time by evaluating
  defects at each edge independently. Locality of the \(\ell^1\) norm
  ensures no global state is required.
\item
  \textbf{Hodge decomposition solver}: Solves the LP
  \(\min_{f \in C^0} \|\delta(s) - d_0 f\|_1\) to separate exact
  (repairable) from harmonic (irreducible) components. Standard network
  simplex or interior-point methods achieve polynomial complexity in the
  graph size.
\item
  \textbf{Defect pattern classifier}: Analyzes the support of the defect
  field \(\delta(s)\) to identify recurring structural patterns (e.g.,
  cycles violating transitivity, conflicting paths, dimensional
  mismatches). This is a topological analysis of the coboundary's
  support structure, implementable via standard graph algorithms.
\item
  \textbf{Epistemic filter}: Given a threshold \(\tau\), accepts outputs
  with \(I(s) \leq \tau\) and rejects outputs with \(I(s) > \tau\). For
  more granular control, applies thresholds to \(R(s)\) (repairable
  defect) and \(\Phi\) (irreducible floor) separately.
\item
  \textbf{Certificate generator}: Constructs the tuple
  \((I, \Phi, R, \eta, \mathcal{C})\) by combining outputs from the
  coboundary auditor, Hodge solver, and cohomology computation. The
  certificate is a structured data object suitable for logging,
  auditing, or contract enforcement.
\end{itemize}

The mathematical framework guarantees that these components compose
correctly: the output of the Hodge decomposition feeds the AGI repair
loop, the repaired section feeds the coboundary auditor for
verification, and the final defect score is provably bounded by
\(\Phi\). No component requires model access or training data --- all
operate on the domain graph and output section.

\begin{center}\rule{0.5\linewidth}{0.5pt}\end{center}

\subsection{12. Conclusion}\label{conclusion}

The \(\ell^1\) coboundary norm is the unique defect measure compatible
with local, additive, faithful observation. This is a theorem, not a
design choice. Current AI systems are trained under \(\ell^2\)
objectives that permit cancellation, interpolation, and hidden
inconsistency. The structural consequence is hallucination --- confident
outputs that correspond to no coherent worldview.

The framework presented here provides:

\begin{enumerate}
\def\labelenumi{\arabic{enumi}.}
\tightlist
\item
  \textbf{A computable hallucination score} \(I(s)\) for any AI output
  on any graph-modeled domain.
\item
  \textbf{A decomposition} into repairable errors (\(R\)) and
  irreducible domain complexity (\(\Phi\)).
\item
  \textbf{A repair algorithm} (AGI) with proved monotone convergence and
  a cohomological floor.
\item
  \textbf{A certificate} \((I, \Phi, R, \eta, \mathcal{C})\) that can be
  audited, composed, and placed in contracts.
\item
  \textbf{A categorical framework} for transferring verification across
  related domains.
\end{enumerate}

The mathematical foundations are proved (Papers 000--005). The
mathematical framework admits efficient implementation. The engineering
path from theorem to deployment is short: build a domain graph, compute
the coboundary, solve the LP, issue the certificate.

AI systems equipped with \(\ell^1\) verification layers admit explicit,
computable bounds on structural inconsistency.

\begin{center}\rule{0.5\linewidth}{0.5pt}\end{center}

\subsection{Appendix: Background for ML
Researchers}\label{appendix-background-for-ml-researchers}

{[}Layer A --- pedagogical exposition of standard mathematical concepts
in ML-accessible terms{]}

This appendix presents the core mathematical concepts (presheaves,
coboundaries, Hodge decomposition) in concrete graph-theoretic language,
without category-theoretic formalism. The goal is to make the framework
accessible to machine learning researchers with undergraduate linear
algebra and graph theory.

\subsubsection{A.1 Domain Graphs as Directed Graphs with
Constraints}\label{a.1-domain-graphs-as-directed-graphs-with-constraints}

\textbf{What is a domain graph?}

A domain graph is just a finite directed graph \(G = (V, E)\) where: -
Each \textbf{vertex} \(v \in V\) represents a concept or claim (e.g.,
``the patient has a fever'', ``temperature = 38.5°C''). - Each
\textbf{edge} \(e = (u, v) \in E\) represents a constraint or
relationship (e.g., ``if Flu then Fever'', ``temperature \textgreater{}
37°C implies fever'').

Unlike standard graph embeddings, we don't learn continuous
representations. Instead, we work with explicit \textbf{sections}:
assignments of values to vertices.

\textbf{What is a section?}

A section \(s\) is a function \(s: V \to F\) that assigns a value
\(s(v) \in F\) to each vertex. The fiber \(F\) is the value space --- it
could be: - \(F = \mathbb{R}\) (real numbers) - \(F = [0, 1]\)
(probabilities) - \(F = \{0, 1\}\) (binary labels) -
\(F = \mathbb{R}^k\) (vectors)

An AI system's output on a domain is a section: it assigns a value
(claim, probability, label) to each concept in the graph.

\textbf{What is a restriction map?}

Each edge \(e = (u, v)\) carries a \textbf{restriction map}
\(r_e: F \to F\) that says: ``if the value at \(v\) is \(x\), what does
that imply for the value at \(u\)?''

Example: If edge \(e\) represents ``if Flu (prob \(p\)) then Fever (prob
\(\geq 0.8p\))'', the restriction map is \(r_e(p) = 0.8p\).

\subsubsection{A.2 The Coboundary: Measuring
Inconsistency}\label{a.2-the-coboundary-measuring-inconsistency}

\textbf{Definition}: The coboundary \(\delta_e(s)\) at edge
\(e = (u, v)\) is:

\[\delta_e(s) = r_e(s(v)) - s(u)\]

This measures the \textbf{mismatch} between: - What vertex \(v\) claims
(projected through the constraint \(e\)): \(r_e(s(v))\) - What vertex
\(u\) claims directly: \(s(u)\)

If \(\delta_e(s) = 0\), the claims are consistent. If
\(\delta_e(s) \neq 0\), there's a contradiction.

\textbf{The total defect} is:

\[I(s) = \sum_{e \in E} |\delta_e(s)|\]

This is just the sum of absolute violations across all edges --- the
\(\ell^1\) norm of the coboundary vector.

\textbf{Why \(\ell^1\) instead of \(\ell^2\)?}

The \(\ell^2\) norm \(\sqrt{\sum_e \delta_e^2}\) permits cancellation:
errors of opposite sign partially cancel. The \(\ell^1\) norm
\(\sum_e |\delta_e|\) does not. Every nonzero defect contributes to the
total. This prevents averaging between contradictory claims.

\subsubsection{A.3 The Hodge Decomposition: Repairable vs.~Irreducible
Defects}\label{a.3-the-hodge-decomposition-repairable-vs.-irreducible-defects}

\textbf{Question}: Given a defect field
\(\delta = (\delta_{e_1}, \delta_{e_2}, \ldots, \delta_{e_m}) \in \mathbb{R}^m\)
(one number per edge), how much of it can be fixed by adjusting the
section \(s\)?

\textbf{Answer}: Solve the following \textbf{linear program}:

\[\Phi = \min_{f: V \to \mathbb{R}} \sum_{e = (u,v) \in E} |\delta_e - (f(v) - f(u))|\]

Here, \(f\) is a \textbf{gauge transformation} --- a function that
shifts the section values. The term \(f(v) - f(u)\) is the coboundary of
\(f\) at edge \(e\), written \(d_0 f(e)\).

\textbf{Interpretation}: - The LP finds the optimal shift \(f\) that
brings \(\delta\) as close as possible to the \textbf{exact form}
\(d_0 f\) (a coboundary of a 0-cochain). - The residual \(\Phi\) is what
remains after the best possible gauge shift. This is the
\textbf{irreducible defect} --- it cannot be removed by changing \(s\)
alone.

\textbf{Decomposition}:

\[\delta = d_0 f^* + r^*\]

where: - \(d_0 f^* = (f^*(v) - f^*(u))\) is the \textbf{exact
(repairable) component} --- fix this by adjusting
\(s(v) \gets s(v) - f^*(v)\) for all \(v\). - \(r^*\) is the
\textbf{harmonic (irreducible) component} with \(\|r^*\|_1 = \Phi\) ---
this cannot be removed without changing the domain graph.

\textbf{Why this matters}:

If \(\Phi > 0\), the domain graph itself contains contradictions. Even a
perfect AI cannot achieve \(I = 0\) on this domain. Knowing \(\Phi\)
tells you where the domain model needs fixing.

\subsubsection{A.4 The Autopoietic Repair
Algorithm}\label{a.4-the-autopoietic-repair-algorithm}

The AGI (Autopoietic Gauge Iterator) repairs the section by iteratively
applying a simple update rule:

\textbf{For each vertex \(v\) with defect}: 1. Collect all neighbors'
projected values: \(\{r_e(s(u)) : e = (u, v) \text{ or } e = (v, u)\}\)
2. Replace \(s(v)\) with the \textbf{\(\ell^1\) median} of these
projected values.

The \textbf{median} (not mean!) ensures no averaging. The update snaps
to a value that minimizes the total \(\ell^1\) deviation from neighbors.

\textbf{Convergence}: Each iteration strictly decreases the total defect
\(I(s)\) until it reaches \(\Phi\) (the irreducible floor). The
algorithm stops when no further improvement is possible.

\textbf{Complexity}: For a graph with \(|V|\) vertices and \(|E|\)
edges, the algorithm converges in \(O(|V| \cdot I_0 / \varepsilon)\)
iterations in the worst case, where \(I_0\) is the initial defect and
\(\varepsilon\) is the tolerance. Under non-degeneracy, convergence is
geometric: \(O(\log(1/\varepsilon))\) iterations.

\subsubsection{A.5 Linear Programming Formulation
(Concrete)}\label{a.5-linear-programming-formulation-concrete}

For readers familiar with LP solvers, here is the explicit formulation
for the Hodge decomposition:

\textbf{Variables}: \(f_1, \ldots, f_n \in \mathbb{R}\) (one per vertex
in \(V = \{v_1, \ldots, v_n\}\))

\textbf{Auxiliary variables}: \(z_1, \ldots, z_m \geq 0\) (one per edge
in \(E\))

\textbf{Objective}:

\[\min \sum_{j=1}^m z_j\]

\textbf{Constraints}: For each edge \(e_j = (v_i, v_k)\):

\[-z_j \leq \delta_{e_j} - (f_k - f_i) \leq z_j\]

This is a standard LP with \(n + m\) variables and \(2m\) inequality
constraints. Any off-the-shelf LP solver (GLPK, HiGHS, Gurobi, CPLEX)
can solve it in polynomial time.

\textbf{Output}: The optimal \(f^* = (f_1^*, \ldots, f_n^*)\) and the
minimal cost \(\Phi = \sum_j z_j^*\).

\subsubsection{A.6 Comparison with Familiar ML
Concepts}\label{a.6-comparison-with-familiar-ml-concepts}

\begin{longtable}[]{@{}
  >{\raggedright\arraybackslash}p{(\linewidth - 2\tabcolsep) * \real{0.3077}}
  >{\raggedright\arraybackslash}p{(\linewidth - 2\tabcolsep) * \real{0.6923}}@{}}
\toprule\noalign{}
\begin{minipage}[b]{\linewidth}\raggedright
ML Concept
\end{minipage} & \begin{minipage}[b]{\linewidth}\raggedright
\(\ell^1\) Framework Analog
\end{minipage} \\
\midrule\noalign{}
\endhead
\bottomrule\noalign{}
\endlastfoot
\textbf{Feature vector} & Section \(s\) (assigns values to vertices) \\
\textbf{Loss function} & Coboundary defect
\(I(s) = \sum_e \|\delta_e\|\) \\
\textbf{Training data} & Domain graph \(G\) (constraints, not
samples) \\
\textbf{Regularization} & Irreducible defect \(\Phi\) (structural lower
bound) \\
\textbf{Gradient descent} & AGI median updates (iterative repair) \\
\textbf{Validation error} & Total defect \(I(s)\) on test graph \\
\textbf{Overfitting} & Section fits one subgraph but violates global
constraints \\
\end{longtable}

\textbf{Key difference}: Standard ML minimizes loss over training
samples. The \(\ell^1\) framework minimizes defect over
\textbf{structural constraints}. It is not statistical --- it is
topological.

\subsubsection{A.7 When to Use This
Framework}\label{a.7-when-to-use-this-framework}

\textbf{Use the \(\ell^1\) coboundary framework when}: - The domain has
explicit constraints (logical, causal, ontological). - Inconsistency has
consequences (medical, legal, financial, safety-critical). - You need a
\textbf{certificate} of consistency, not just a probability. - The
deployment domain differs from the training domain
(out-of-distribution).

\textbf{Do not use this framework when}: - The domain has no explicit
constraints (purely perceptual tasks like image recognition). -
Statistical uncertainty is the primary concern (use conformal
prediction). - You only care about average-case performance (use
standard \(\ell^2\) metrics).

\subsubsection{A.8 Minimal Working Example (Python
Pseudocode)}\label{a.8-minimal-working-example-python-pseudocode}

\begin{Shaded}
\begin{Highlighting}[]
\ImportTok{import}\NormalTok{ numpy }\ImportTok{as}\NormalTok{ np}
\ImportTok{from}\NormalTok{ scipy.optimize }\ImportTok{import}\NormalTok{ linprog}

\CommentTok{\# Define domain graph}
\NormalTok{V }\OperatorTok{=}\NormalTok{ [}\StringTok{\textquotesingle{}Fever\textquotesingle{}}\NormalTok{, }\StringTok{\textquotesingle{}Headache\textquotesingle{}}\NormalTok{, }\StringTok{\textquotesingle{}Flu\textquotesingle{}}\NormalTok{, }\StringTok{\textquotesingle{}Migraine\textquotesingle{}}\NormalTok{, }\StringTok{\textquotesingle{}Dehydration\textquotesingle{}}\NormalTok{]}
\NormalTok{E }\OperatorTok{=}\NormalTok{ [(}\StringTok{\textquotesingle{}Flu\textquotesingle{}}\NormalTok{, }\StringTok{\textquotesingle{}Fever\textquotesingle{}}\NormalTok{, }\FloatTok{0.8}\NormalTok{), (}\StringTok{\textquotesingle{}Flu\textquotesingle{}}\NormalTok{, }\StringTok{\textquotesingle{}Headache\textquotesingle{}}\NormalTok{, }\FloatTok{0.6}\NormalTok{), ...]  }\CommentTok{\# (u, v, weight)}

\CommentTok{\# AI output (section)}
\NormalTok{s }\OperatorTok{=}\NormalTok{ \{}\StringTok{\textquotesingle{}Fever\textquotesingle{}}\NormalTok{: }\FloatTok{0.7}\NormalTok{, }\StringTok{\textquotesingle{}Headache\textquotesingle{}}\NormalTok{: }\FloatTok{0.8}\NormalTok{, }\StringTok{\textquotesingle{}Flu\textquotesingle{}}\NormalTok{: }\FloatTok{0.6}\NormalTok{, }\StringTok{\textquotesingle{}Migraine\textquotesingle{}}\NormalTok{: }\FloatTok{0.5}\NormalTok{, }\StringTok{\textquotesingle{}Dehydration\textquotesingle{}}\NormalTok{: }\FloatTok{0.4}\NormalTok{\}}

\CommentTok{\# Compute coboundary defects}
\NormalTok{defects }\OperatorTok{=}\NormalTok{ []}
\ControlFlowTok{for}\NormalTok{ (u, v, w) }\KeywordTok{in}\NormalTok{ E:}
\NormalTok{    delta\_e }\OperatorTok{=}\NormalTok{ w }\OperatorTok{*}\NormalTok{ s[v] }\OperatorTok{{-}}\NormalTok{ s[u]  }\CommentTok{\# simplified restriction map: r\_e(x) = w*x}
\NormalTok{    defects.append(}\BuiltInTok{abs}\NormalTok{(delta\_e))}

\CommentTok{\# Total defect}
\NormalTok{I }\OperatorTok{=} \BuiltInTok{sum}\NormalTok{(defects)}
\BuiltInTok{print}\NormalTok{(}\SpecialStringTok{f"Total defect: }\SpecialCharTok{\{}\NormalTok{I}\SpecialCharTok{\}}\SpecialStringTok{"}\NormalTok{)}

\CommentTok{\# Hodge decomposition (LP formulation)}
\CommentTok{\# Variables: f[v] for each vertex v}
\CommentTok{\# Minimize: sum |delta\_e {-} (f[v] {-} f[u])|}
\NormalTok{c }\OperatorTok{=}\NormalTok{ np.ones(}\BuiltInTok{len}\NormalTok{(E))  }\CommentTok{\# objective: minimize sum of z\_j}
\CommentTok{\# ... construct A\_ub, b\_ub from constraints {-}z\_j \textless{}= delta\_e {-} (f\_v {-} f\_u) \textless{}= z\_j}
\CommentTok{\# res = linprog(c, A\_ub, b\_ub)}
\CommentTok{\# Phi = res.fun}

\CommentTok{\# AGI repair (iterative median updates)}
\ControlFlowTok{for}\NormalTok{ iteration }\KeywordTok{in} \BuiltInTok{range}\NormalTok{(max\_iters):}
    \ControlFlowTok{for}\NormalTok{ v }\KeywordTok{in}\NormalTok{ V:}
\NormalTok{        neighbors }\OperatorTok{=}\NormalTok{ [...]  }\CommentTok{\# collect projected values from neighbors}
\NormalTok{        s[v] }\OperatorTok{=}\NormalTok{ np.median(neighbors)  }\CommentTok{\# L1 median update}
\NormalTok{    I\_new }\OperatorTok{=}\NormalTok{ compute\_total\_defect(s)}
    \ControlFlowTok{if}\NormalTok{ I\_new }\OperatorTok{\textgreater{}=}\NormalTok{ I }\OperatorTok{{-}}\NormalTok{ epsilon:}
        \ControlFlowTok{break}  \CommentTok{\# converged}
\NormalTok{    I }\OperatorTok{=}\NormalTok{ I\_new}

\CommentTok{\# Certificate}
\NormalTok{certificate }\OperatorTok{=}\NormalTok{ (I, Phi, I }\OperatorTok{{-}}\NormalTok{ Phi, (I }\OperatorTok{{-}}\NormalTok{ Phi) }\OperatorTok{/}\NormalTok{ I)}
\BuiltInTok{print}\NormalTok{(}\SpecialStringTok{f"Certificate: }\SpecialCharTok{\{}\NormalTok{certificate}\SpecialCharTok{\}}\SpecialStringTok{"}\NormalTok{)}
\end{Highlighting}
\end{Shaded}

This pseudo code captures the entire pipeline: domain graph
construction, coboundary computation, Hodge decomposition via LP, AGI
repair, and certificate generation.

\begin{center}\rule{0.5\linewidth}{0.5pt}\end{center}

\subsection{References}\label{references}

{[}1{]} J. H. Carroll, ``Projection Obstruction Theory: Retraction
Nonlinearity, \(\ell^1\) Rigidity, and Density Scaling,'' 2026. DOI:
10.5281/zenodo.19151803

{[}2{]} J. H. Carroll, ``The Free \(\ell^1\) Seminorm on Banach Presheaf
Coboundaries,'' 2026. DOI: 10.5281/zenodo.19152115

{[}3{]} J. H. Carroll, ``Coordinate-Wise Additivity and the \(\ell^1\)
Norm on Finite Graph Cochains,'' 2026. DOI: 10.5281/zenodo.19152599

{[}4{]} J. H. Carroll, ``Hodge Structure and Gauge Equivalence of
\(\ell^1\) Defect Fields,'' 2026. DOI: 10.5281/zenodo.19152799

{[}5{]} J. H. Carroll, ``Autopoietic Cohomology: Iterative Obstruction
Repair, Manifold Emergence, and Gauge Structure on Causal Posets,''
2026. DOI: 10.5281/zenodo.19155011

{[}6{]} J. H. Carroll, ``A Common \(\ell^1\) Defect Framework for
Quantum, Causal, and Gravitational Systems,'' 2026. DOI:
10.5281/zenodo.19155079

{[}7{]} J. H. Carroll, ``Constraint Closure and the \(\ell^1\) Structure
of Physical Law,'' 2026. DOI: 10.5281/zenodo.19183020

{[}8{]} A. Chambolle, ``An Algorithm for Total Variation Minimization
and Applications,'' \emph{J. Math. Imaging Vis.} 20, 89--97, 2004.

{[}9{]} E. J. Candes and T. Tao, ``Decoding by Linear Programming,''
\emph{IEEE Trans. Inf. Theory} 51(12), 4203--4215, 2005.

{[}10{]} L. Ambrosio, N. Gigli, and G. Savare, \emph{Gradient Flows in
Metric Spaces and in the Space of Probability Measures}, Birkhauser,
2008.

{[}11{]} A. Vaswani et al., ``Attention Is All You Need,''
\emph{NeurIPS}, 2017.

{[}12{]} Y. Bengio, A. Courville, and P. Vincent, ``Representation
Learning: A Review and New Perspectives,'' \emph{IEEE Trans. PAMI}
35(8), 2013.

\end{document}
